\documentclass[12pt,openany,a4paper]{book}
\usepackage{graphics}
\usepackage{array,makecell,graphicx,pifont,tabularx,amssymb}
\usepackage{rotating}
\usepackage{lscape}
\usepackage{multirow}
\usepackage{adjustbox}
\usepackage[a4paper,margin=1in]{geometry}
\usepackage[dvipsnames]{xcolor}
\usepackage{enumitem}
\usepackage{booktabs}
\usepackage{amsmath}
\usepackage{hyperref}

\usepackage{titlesec}
\titleformat{\chapter}[hang]
  {\normalfont\huge\bfseries}
  {\thechapter.\ }{0pt}{}
\titlespacing*{\chapter}{0pt}{0pt}{1em}

\setcounter{secnumdepth}{2}
\setcounter{tocdepth}{2}

\pagestyle{headings}
\raggedbottom

\setlength{\topmargin}     {-5mm}
\setlength{\oddsidemargin} {10mm}
\setlength{\evensidemargin}{0mm}
\setlength{\textwidth}     {150mm}
\setlength{\textheight}    {240mm}

\renewcommand{\baselinestretch}{1.2}

\renewcommand{\textfraction}{0.25}
\renewcommand{\topfraction}{0.75}
\renewcommand{\bottomfraction}{0.75}
\renewcommand{\floatpagefraction}{0.75}

\newtheorem{defn}   {Definition}
\newtheorem{theorem}{Theorem}
\newtheorem{lemma}  {Lemma}

\newcommand{\eref}[1]{(\ref{#1})}
\newcommand{\eq}[1]  {Eq.\,(\ref{#1})}
\newcommand{\fig}[1] {Fig.\,\ref{#1}}
\newcommand{\tab}[1] {Table~\ref{#1}}
\newcommand{\secn}[1]{Section~\ref{#1}}

\newcommand{\teq}[1]{\mbox{$#1$}}
\newcommand{\qed}{\hspace*{\fill}$\bullet$}
\newcommand{\cmark}{\ding{51}}

% Research-specific macros
\newcommand{\taug}{\tau_{\mathrm{gen}}}
\newcommand{\tauf}{\tau_{F}}
\newcommand{\Dtau}{\Delta\tau}
\newcommand{\fcorr}{F_{\mathrm{corr}}}

\begin{document}

\frontmatter

% ================================================================
%  Title Page
% ================================================================
\begin{titlepage}
\renewcommand{\baselinestretch}{1.0}
\begin{center}
\vspace*{35mm}
\Huge\bf

        
\includegraphics[width=60mm]{uq_logo.png}\\
        \vspace*{25mm}

        Does Generalisation Precede Fourier Alignment\\
        in Transformer Grokking?\\
        A Study of the Hyperparameter Phase Diagram\\

\vspace{20mm}
\large\sl
        by\\
        Yueran Ma
        \medskip\\

\large\sl
        under the supervision of\\
        Marcus Gallagher
        \medskip\\
\rm

\end{center}
\end{titlepage}

% ================================================================
%  Use of AI Statement
% ================================================================
\chapter{Use of AI Statement}

\noindent \textit{I acknowledge the use of generative AI tools in completing this
assessment. Details of which tools were used and how they were used are
summarised in Table~\ref{tab:ai_usage}, along with appropriate in-text and full
references. I take responsibility for critically evaluating and
integrating the AI-generated content, and ensuring it adheres to
academic integrity standards.}

\section*{Have you used AI to}
\begin{itemize}[leftmargin=1.5em]
  \item \textbf{Explore your topic?} Yes. Claude Code and ChatGPT were used to
    explore prior literature on grokking, Fourier circuits, and representation
    learning dynamics, providing summaries of key papers and suggesting
    relevant related work.
  \item \textbf{Define your research question?} No. The research question was
    defined independently; AI tools were not used for this purpose.
  \item \textbf{Assist with your literature review?} Yes. AI tools were used to
    summarise and compare key papers (Power et al.\ 2022; Nanda et al.\ 2023;
    Liu et al.\ 2022), though all claims were independently verified against
    original sources.
  \item \textbf{Assist with formatting?} Yes. Grammar checking and LaTeX
    formatting assistance were provided by AI tools.
  \item \textbf{Other aspects of your work?} Yes. Claude Code assisted with
    Python code generation for experiment scripts, metric implementations,
    and result visualisation. All generated code was reviewed and tested
    before use.
\end{itemize}

\begin{table}[htbp]
\centering
\caption{AI tool usage (tick all that apply).}
\label{tab:ai_usage}
\begin{tabular}{|c|c|c|c|c|c|c|c|c|c|}
\hline
\multirow{2}{*}{\parbox{2.5cm}{\vspace*{\fill}\centering AI Model and date\vspace*{0pt}}}
 & \adjustbox{rotate=90}{Language Translation}
 & \adjustbox{rotate=90}{Grammar / style / spelling}
 & \adjustbox{rotate=90}{Topic exploration}
 & \adjustbox{rotate=90}{Research question}
 & \adjustbox{rotate=90}{Literature review}
 & \adjustbox{rotate=90}{Content creation visual (formatting)}
 & \adjustbox{rotate=90}{Content creation (e.g.\ code, text)}
 & \adjustbox{rotate=90}{Feedback}
 & \adjustbox{rotate=90}{Other (provide details)} \\ \hline
\makecell{ChatGPT 4o\\Weeks 1--4}
 & & \cmark & \cmark & & \cmark & & \cmark & & \\ \hline
\makecell{Claude Code\\Weeks 1--12}
 & & \cmark & \cmark & & \cmark & & \cmark & & \\ \hline
 & & & & & & & & & \\ \hline
 & & & & & & & & & \\ \hline
\end{tabular}
\end{table}

% ================================================================
%  Abstract
% ================================================================
\chapter{Abstract}

This project examines how two milestones unfold over time during transformer grokking on modular arithmetic:
the onset of test generalisation~$\taug$ and the onset of
Fourier-aligned token embedding geometry~$\tauf$.
This distinction matters because the Fourier circuit identified by Nanda et al.\ \cite{nanda2023progress} is the known mechanism behind grokked solutions, but it is still unclear whether the \emph{embedding-level} Fourier geometry settles into place before or after the model begins to generalise.
To answer this, we will train a decoder-only transformer on $(a+b)\bmod 53$ across fine-grained sweeps of decoder weight decay and learning rate,
measuring $\taug$ via a sustained threshold detector and $\tauf$ via a
BIC changepoint on permutation-null-corrected Fourier alignment.
By sweeping both axes of the hyperparameter phase diagram, we aim to test whether G$<$F (generalisation preceding embedding Fourier alignment) appears as a robust and repeatable pattern within the Grokking phase, and to quantify the signed lag $\Dtau = \tauf - \taug$ across hyperparameter settings.
A sensitivity analysis across multiple detector configurations will then check how robust that ordering is.

\tableofcontents

\listoffigures
\addcontentsline{toc}{chapter}{List of Figures}

\listoftables
\addcontentsline{toc}{chapter}{List of Tables}

\cleardoublepage
\mainmatter

% ================================================================
\chapter{Introduction}
% ================================================================

Grokking---the abrupt onset of test generalisation long after training
accuracy saturates---has emerged as a canonical testbed for studying
delayed representation learning in neural networks
\cite{power2022grokking}.
Unlike standard learning dynamics, grokking exhibits a sharp transition
from near-zero test accuracy to near-perfect test accuracy after
thousands of additional training steps, despite the training loss
having converged much earlier.
Understanding \emph{why} generalisation is delayed, and what internal
changes accompany the transition, is a central open problem in the
study of neural network learning dynamics.

On modular arithmetic tasks, mechanistic interpretability
\cite{olah2020circuits} has provided a compelling explanation:
grokked transformers implement a Fourier circuit
\cite{nanda2023progress}.
The network computes $(a+b)\bmod p$ by exploiting trigonometric
identities over a small set of task-relevant harmonics.
One striking geometric consequence of this algorithm is that the $p$ token embeddings form a near-circular ring in PCA space, with embedding $i$ positioned at angle $2\pi i/p$---an arrangement that has become one of the most recognisable images in the grokking literature.

Nanda et al.\ \cite{nanda2023progress} further showed that mechanism-level
\emph{progress measures}---quantities such as the Fourier-squared
loss---improve during a circuit-formation phase that \emph{precedes}
the generalisation jump.
This naturally leads to a closely related question about token \emph{embedding geometry}: does embedding-level Fourier alignment also settle in before generalisation, or does it arrive later?

\section{Research Question}

This project asks when two events occur relative to one another along the training trajectory of a grokking transformer:
\begin{itemize}
  \item $\taug$: generalisation onset---the first checkpoint where test
    accuracy $\geq 0.9$ and remains so for three consecutive evaluations;
  \item $\tauf$: embedding-geometry onset---the BIC-selected changepoint
    in the null-corrected Fourier alignment score $\fcorr(t)$, indicating
    the first sustained rise of token embedding Fourier structure above
    a permutation-null baseline.
\end{itemize}
The primary outcome variable is the signed lag $\Dtau = \tauf - \taug$.
Positive $\Dtau$ indicates \textbf{G$<$F} (generalisation precedes
embedding Fourier alignment); negative indicates \textbf{F$<$G}.

Prior work has not directly measured $\Dtau$.
Liu et al.\ \cite{liu2022towards} built a static phase diagram but did
not measure within-run event timing.
Nanda et al.\ \cite{nanda2023progress} studied mechanism-level dynamics
but did not separately track embedding geometry.
This project makes two contributions: (i)~a methodological contribution
in the form of a null-calibrated, BIC-based detector for embedding-level
Fourier alignment onset---designed to be robust across hyperparameter
configurations---and (ii)~an empirical contribution in the form of
the first systematic measurement of $\Dtau$ across the hyperparameter
phase diagram, addressing whether G$<$F holds as a reproducible
structural regularity of transformer grokking.
Beyond grokking itself, resolving the order in which generalisation
and embedding-level Fourier structure appear bears on two broader
questions in mechanistic interpretability and late-training dynamics:
whether interpretable geometric structure---of the kind also studied
under the linear representation hypothesis \cite{park2024linear}---is
a \emph{cause} or a \emph{consequence} of generalisation, and whether
late-training
representation consolidation in transformers belongs to the same
family of post-convergence simplification phenomena as neural collapse
\cite{papyan2020prevalence}.

\subsection{Overview}

The rest of this proposal is organised as follows.
Chapter~2 reviews grokking, the Fourier circuit, and the existing phase-diagram literature, then identifies the specific knowledge gap addressed here.
Chapter~3 presents the project plan, including the task and model specification, metric definitions, experimental design, expected outcomes, and timeline.

% ================================================================
\chapter{Background and Literature Review}
% ================================================================

\section{Grokking and the Hyperparameter Phase Diagram}

Grokking was first documented by Power et al.\ \cite{power2022grokking}
in small transformers trained on binary operations including
$(a+b)\bmod p$.
The defining characteristic is a sharp transition to near-perfect test
accuracy that occurs thousands of gradient steps after training
accuracy has saturated, creating a prolonged period of apparent
``memorisation'' before the model generalises.

Liu et al.\ \cite{liu2022towards} extended this picture by mapping
training behaviour across a two-dimensional (decoder learning rate,
weight decay) hyperparameter grid, identifying four qualitatively
distinct phases:
\begin{itemize}
  \item \textbf{Grokking}: training accuracy saturates early; test
    accuracy jumps to near-perfect much later (delayed generalisation).
  \item \textbf{Memorisation}: training accuracy saturates; test
    accuracy never reaches the threshold (no generalisation).
  \item \textbf{Comprehension}: both training and test accuracy rise
    together rapidly (immediate generalisation).
  \item \textbf{Confusion}: neither training nor test accuracy reaches
    the threshold (learning failure).
\end{itemize}
Their effective-theory analysis attributed phase transitions to
competition between memorisation and generalisation mechanisms,
regulated by weight decay acting as a complexity penalty
\cite{krogh1991weightdecay}.

A growing body of recent work has further refined the picture.
Thilak et al.\ \cite{thilak2022slingshot} reported a
\emph{slingshot} mechanism in which adaptive optimisers transiently
amplify the loss before generalisation, identifying optimiser
dynamics as a contributing cause of delayed generalisation.
Notsawo et al.\ \cite{notsawo2023predicting} showed that grokking can
be predicted from early-training spectral signatures, suggesting that
the eventual generalisation jump is foreshadowed in measurable
representation statistics.
Davies et al.\ \cite{davies2023unifying} cast grokking and double
descent as two facets of the same underlying capacity-versus-data
tradeoff, situating it within classical generalisation theory.
Merrill et al.\ \cite{merrill2023tale} proposed an alternative
``two-circuit'' account in which a sparse generalising subnetwork and
a dense memorising subnetwork compete throughout training, offering
a competing explanation that our embedding-level timing measurement
can in principle help discriminate.

While Liu et al.\ \cite{liu2022towards} studied a specific algorithmic setting,
a follow-up paper \cite{liu2023omnigrok} demonstrated that grokking is not
confined to modular arithmetic: it appears across image classification,
sparse parity, and other tasks when weight norms are appropriately
controlled, suggesting that the phase structure is a general property
of regularised learning rather than an artefact of the task.
On the theoretical side, Lyu et al.\ \cite{lyu2024dichotomy} attributed
grokking to a \emph{dichotomy of implicit biases}: early-phase gradient
descent favours low-loss but high-complexity solutions, while late-phase
dynamics favour sparse, generalisable solutions.
This principled account of why the generalisation transition is delayed
reinforces the view that the training trajectory contains discrete
structural stages---precisely the kind of staged dynamics our timing
analysis aims to characterise at the embedding level.

Under the specific configuration used in this project ($p=53$,
$d_{\mathrm{model}}=256$, train fraction~$0.3$), we expect
the Comprehension phase to be absent or rare in the hyperparameter
region to be tested, and anticipate observing primarily Grokking
and Memorisation/Confusion phases.
One possible reason is that the model has excess capacity relative
to the training set size at $p=53$ (approximately 840 training
pairs), though the sweep range, threshold definitions, and task
choice may also contribute to the phase boundaries observed.

\section{The Fourier Circuit and Embedding Geometry}

Nanda et al.\ \cite{nanda2023progress} reverse-engineered the
algorithm implemented by grokked transformers on $(a+b)\bmod p$:
the network computes modular sums via trigonometric identities
over task-relevant Fourier harmonics $K\subset\{1,\dots,\lfloor
p/2\rfloor\}$.
Concretely, the logit for output $c$ is a sum of products of cosines:
\[
  \mathrm{logit}(c) \;\propto\; \sum_{k\in K}
    \cos\!\left(\tfrac{2\pi k}{p}(a+b-c)\right).
\]
As a geometric consequence, the embedding of token $i$ lies near
the point $(\cos(2\pi ki/p),\sin(2\pi ki/p))$ in the space spanned
by the $k$-th harmonic components, giving rise to the characteristic
near-circular ring arrangement observable in PCA projections.

Importantly, Nanda et al.\ also measured \emph{progress measures}---
logit-level and weight-level Fourier statistics---and found that these
measures improve \emph{before} the generalisation jump, during what
they called the ``circuit formation'' phase.
A subsequent ``cleanup'' phase follows, in which memorisation
components are pruned.

Varma et al.\ \cite{varma2023explaining} explained grokking through
\emph{circuit efficiency}: generalising circuits are more weight-efficient
than memorising circuits, so weight decay preferentially prunes
memorisation components during the cleanup phase.
Under this view, the cleanup phase is driven by regularisation acting
on an already-formed generalising circuit---providing a mechanistic
motivation for why embedding-level Fourier geometry (a geometric
signature of the generalising circuit) may only fully consolidate
\emph{after} generalisation onset, as the memorisation components
are removed.

Zhong et al.\ \cite{zhong2023clock} proposed two complementary
mechanistic stories for modular arithmetic: the ``Clock'' algorithm
(consistent with Nanda et al.'s Fourier view) and the ``Pizza''
algorithm (based on an alternative geometric decomposition).
Both predict Fourier-structured token embeddings as a consequence of
learning the task.
Crucially, $\fcorr(t)$ is agnostic to which algorithm the network
implements---it measures the degree of Fourier structure in the
embeddings regardless of the underlying computational mechanism---so
the timing analysis applies under either mechanistic story.

Gromov \cite{gromov2023grokking} provided an algebraic characterisation
of grokking circuits, showing that networks exploit the group structure
of $\mathbb{Z}/p\mathbb{Z}$ and derived closed-form scaling laws for
the grokking delay.
This algebraic perspective confirms that the Fourier basis is the
natural coordinate system for both the circuit (weight space) and
the token embeddings (representation space), motivating $\fcorr(t)$
as a principled embedding-level probe.

This project complements that body of work by measuring a
related but distinct quantity: the geometric organisation of the
\emph{token embeddings themselves}, as characterised by our
null-corrected Fourier alignment score $\fcorr(t)$.

\subsection{Adjudicating Among Mechanistic Explanations}
\label{sec:adjudicate}

The accounts of Nanda et al.\ \cite{nanda2023progress},
Varma et al.\ \cite{varma2023explaining}, and
Lyu et al.\ \cite{lyu2024dichotomy} are mutually compatible at the
high level---all three identify discrete training stages---but they
make different claims about \emph{when} interpretable representation
structure should appear relative to $\taug$.
Nanda et al.\ argue that mechanism-level Fourier statistics rise
during a circuit-formation phase that \emph{precedes} $\taug$;
on a strict reading, the same should hold of the embedding-level
geometry that the circuit reads from, predicting F$<$G.
Varma et al.\ recast the post-$\taug$ ``cleanup'' phase as
weight-decay-driven pruning of memorisation components from an
already-formed generalising circuit; this implies that
\emph{embedding} geometry, which is dominated by memorisation
components until cleanup, should consolidate \emph{after} $\taug$,
predicting G$<$F.
Lyu et al.'s implicit-bias dichotomy is agnostic about embedding
geometry per~se but predicts that any interpretable late-phase
structure should co-emerge with the late-phase sparse solution,
consistent with G$\approx$F or weak G$<$F.
The signed lag $\Dtau$ measured at the embedding level therefore
discriminates among these three accounts in a way that neither the
logit-level progress measures of Nanda et al.\ nor the
end-of-training weight norms of Liu et al.\ do, and constitutes
the methodological case for an embedding-level, time-resolved probe.

\section{Representation Learning Dynamics}

The broader question of when and how internal representations change
during learning is a central topic in deep learning theory.
Saxe et al.\ \cite{saxe2014exact} derived exact solutions to learning
dynamics in deep linear networks, establishing that the singular
value decomposition of the weight matrices evolves in a staged,
sequential manner: different ``modes'' of the solution switch on at
different times.
Although grokking transformers are highly nonlinear, this theoretical
framework established the principle that internal representations
can reorganise in discrete, temporally separated stages---a
phenomenon our project examines empirically at the level of token
embedding geometry.

Barak et al.\ \cite{barak2022hidden} demonstrated that SGD can learn
parity functions near the computational limit, showing that abrupt
generalisation transitions can hide gradual progress in standard
loss-scale metrics.
This timing analysis aims to test whether token geometry consolidates
after generalisation---a structural change that would be invisible
in loss or accuracy curves alone.

Papyan et al.\ \cite{papyan2020prevalence} documented \emph{neural
collapse}---a geometric simplification of penultimate-layer
representations that emerges \emph{after} training accuracy saturates,
during what they term the terminal phase of training.
The analogy to our setting is direct: just as neural collapse is a
post-convergence consolidation of internal geometry in classification
networks, G$<$F (if confirmed) would identify a post-generalisation
consolidation of embedding geometry in grokking transformers,
situating the phenomenon within a broader pattern of late-training
representation simplification.

Elhage et al.\ \cite{elhage2021mathematical} introduced the
transformer circuits framework, providing mechanistic vocabulary
for understanding how attention and MLP layers implement specific
computations, which underpins the Fourier circuit analysis of
Nanda et al.

Rubin et al.\ \cite{rubin2024grokking} characterised grokking as
a first-order phase transition in two-layer networks, providing a
theoretical basis for the sharp generalisation onset observed
empirically.

\section{Knowledge Gap}

Despite substantial progress in understanding \emph{when} generalisation
occurs and \emph{what mechanism} produces it, the relative timing of
two distinct milestones has not been directly measured:
(i)~the onset of test generalisation $\taug$, and
(ii)~the onset of Fourier-structured token embedding geometry $\tauf$.

Specifically:
\begin{itemize}
  \item Nanda et al.\ \cite{nanda2023progress} showed that circuit
    formation (mechanism-level) precedes $\taug$, but they did not
    separately track embedding geometry onset.
  \item Liu et al.\ \cite{liu2022towards} constructed the phase diagram
    but measured only static (end-of-training) representations, not
    within-run event timing.
  \item No prior study has systematically measured $\Dtau = \tauf - \taug$
    across the hyperparameter phase diagram.
\end{itemize}

This project differs from those works along three concrete axes.
First, the unit of analysis is the \emph{token embedding matrix} $E_t$,
rather than Nanda et al.'s logit-level Fourier statistic or Liu et al.'s
end-of-training weight norms; this isolates representation geometry
from downstream readout layers.
Second, the metric is \emph{null-corrected} via per-step permutation
baselines (\eq{eq:f_raw}), which removes the high chance-level
alignment present when $d \gg p$ and which neither prior work
accounts for at the embedding level.
Third, the analysis is \emph{within-run and time-resolved}: $\taug$
and $\tauf$ are extracted from the same trajectory via a sustained
threshold and a BIC changepoint, so the signed lag $\Dtau$ is
directly measurable on a per-run basis rather than inferred from
aggregated curves.
Together these choices yield a per-cell, per-seed measurement of
$\Dtau$ that the existing literature does not provide.

Answering whether G$<$F or F$<$G holds---and whether the answer is
consistent across hyperparameter configurations---has direct implications
for our understanding of the relationship between algorithm learning
and representation consolidation in grokking.

% ================================================================
\chapter{Project Plan}
\label{chap:plan}
% ================================================================

\section{Project Aims}

This project aims to:
\begin{enumerate}
  \item Develop a robust, null-calibrated metric for detecting the
    onset of Fourier-aligned token embedding geometry ($\tauf$) during
    transformer training.
  \item Measure the signed lag $\Dtau = \tauf - \taug$ systematically
    across the hyperparameter phase diagram (weight decay axis and
    learning rate axis).
  \item Determine whether G$<$F (generalisation precedes Fourier
    embedding alignment) holds as a consistent, robust pattern within
    the Grokking phase.
  \item Assess the robustness of any detected G$<$F ordering to
    detector hyperparameter choices via a sensitivity analysis.
\end{enumerate}

\section{Methodology}

\subsection{Task and Model}

We will train a decoder-only transformer \cite{vaswani2017attention},
following the GPT-style autoregressive architecture
\cite{radford2018gpt}, on $(a+b)\bmod 53$.
All $p^2 = 2{,}809$ input pairs will be generated; 30\% (843 pairs)
will form the training set, randomly sampled.
The model architecture follows Liu et al.\ \cite{liu2022towards}:
$d_{\mathrm{model}}=256$, 4 attention heads, 2 transformer layers,
feed-forward dimension 1024, no dropout.
Standard architectural components---residual connections
\cite{he2016resnet} and layer normalisation \cite{ba2016layernorm}---are
used as in the reference implementation.
The choice of $p=53$, $d_{\mathrm{model}}=256$, and train fraction
$0.3$ is deliberate: $p=53$ is large enough that the chance-level
Fourier alignment in \eq{eq:f_raw} is non-trivial (so the null
correction is non-vacuous) yet small enough that a single training
run, including the per-checkpoint 100-permutation null estimation,
fits comfortably in 12\,GB of GPU memory; $d=256$ matches the
Liu et al.\ phase diagram so our hyperparameter coordinates are
directly comparable to theirs; and the $0.3$ train fraction sits
inside the regime in which Liu et al.\ report a clear Grokking phase
for $(a+b)\bmod p$, maximising the number of cells in which $\taug$
and $\tauf$ are both expected to be detectable.
The embedding layer will have fixed learning rate $\eta_{\mathrm{emb}}
= 10^{-3}$ and weight decay~0; only decoder parameters will be swept.
Optimiser: AdamW \cite{loshchilov2019adamw} (the decoupled
weight-decay variant of Adam \cite{kingma2015adam}) with 10-step
linear warmup, constant learning rate thereafter, and gradient
clipping at norm~1.0.
Training will run to 50{,}000 steps by default, extended to 80{,}000
if training accuracy exceeds~0.9 but $\taug$ has not been detected.

\subsection{Metrics}

\paragraph{Generalisation onset $\taug$.}
Test accuracy will be evaluated at fixed 500-step intervals.
$\taug$ will be defined as the first checkpoint where test accuracy
$\geq 0.9$ for three consecutive evaluations (1{,}500 steps of
sustained performance above threshold).
Requiring three consecutive evaluations reduces the chance of mistaking a brief accuracy spike for true generalisation onset.

\paragraph{Raw Fourier alignment $F_{\mathrm{raw}}(t)$.}
Let $E_t \in \mathbb{R}^{p \times d}$ be the mean-row-centred token
embedding matrix at step $t$.
The raw Fourier alignment is:
\begin{equation}
  F_{\mathrm{raw}}(t) \;=\;
  \max_{k=1}^{\lfloor p/2\rfloor}
  \!\left(1 - \frac{\|E_t - Q_k Q_k^\top E_t\|_F^2}
  {\|E_t\|_F^2}\right),
  \label{eq:f_raw}
\end{equation}
where $Q_k$ is an orthonormal basis for the subspace spanned by
$[\cos(2\pi k i/p),\sin(2\pi k i/p)]_{i=0}^{p-1}$.
This quantity measures the fraction of embedding variance explained
by the best-fitting single Fourier harmonic (a lower bound on
alignment with the full set of task-relevant harmonics).

\paragraph{Null-corrected Fourier alignment $\fcorr(t)$.}
With $p=53$ and $d=256$, any embedding matrix achieves non-trivial
$F_{\mathrm{raw}}$ by chance due to high dimensionality.
Every 5{,}000 steps we will estimate the 95th percentile of
$F_{\mathrm{raw}}$ under 100 random token-index permutations of the
current embedding matrix:
\begin{equation}
  F_{\mathrm{null,p95}}(t) \;=\;
  \mathrm{Pct}_{95}\!\left(
    \{F_{\mathrm{raw}}(E_t[\sigma]) \mid \sigma \in \mathcal{S}^{(100)}\}
  \right),
\end{equation}
where $\mathcal{S}^{(100)}$ denotes a set of 100 i.i.d.\ uniform random permutations of $\{0,\dots,p{-}1\}$,
and set $\fcorr(t) = \max(0,\, F_{\mathrm{raw}}(t) - F_{\mathrm{null,p95}}(t))$.
Between update steps, $F_{\mathrm{null,p95}}$ will be held constant at its
most recently estimated value.
This removes chance-level baseline alignment and marks the step at which the embeddings show above-chance Fourier structure relative to the empirical permutation null. It is not a formal per-step significance test, but it provides a practical threshold for changepoint detection.
Figure~\ref{fig:raw_vs_corr} compares $F_{\mathrm{raw}}$ and $\fcorr$ across pilot runs and motivates why the null correction is necessary in this regime.

\begin{figure}[htbp]
  \centering
  \includegraphics[width=0.85\textwidth]{raw_vs_corr_scatter.pdf}
  \caption[Raw versus null-corrected Fourier alignment.]{Raw versus null-corrected Fourier alignment across pilot runs. The null correction removes the high baseline that arises from $d \gg p$ and isolates the above-chance component used for changepoint detection.}
  \label{fig:raw_vs_corr}
\end{figure}

\paragraph{Fourier geometry onset $\tauf$.}
We will apply a BIC-selected
\cite{schwarz1978bic} 1-breakpoint segmented regression
\cite{muggeo2003segmented} to $\fcorr(t)$ on a log-time scale
after exponential-moving-average smoothing ($\alpha = 0.15$).
A changepoint will be accepted only if the 4-parameter model strictly
improves BIC over the 2-parameter no-break baseline and the
post-break slope exceeds 1\% of the total value range.
$\tauf$ will be the accepted breakpoint step; labelled ``not detected''
otherwise.
Figure~\ref{fig:timeline} illustrates a representative run with both $\taug$ and $\tauf$ marked on the same time axis.

\begin{figure}[htbp]
  \centering
  \includegraphics[width=0.85\textwidth]{timeline_example.pdf}
  \caption[Representative training timeline.]{Representative timeline of test accuracy and null-corrected Fourier alignment $\fcorr(t)$ for a Grokking-phase run. The vertical lines mark the detected $\taug$ (sustained test accuracy threshold) and $\tauf$ (BIC-selected changepoint), with the signed lag $\Dtau = \tauf - \taug$ as the primary outcome.}
  \label{fig:timeline}
\end{figure}

\paragraph{Phase and ordering labels.}
A run is \textit{Grokking} if both $\tau_{\mathrm{train}}$
(train accuracy $\geq0.9$ for three consecutive evaluations)
and $\taug$ are detected and $\taug - \tau_{\mathrm{train}}
\geq 2{,}000$ steps.
\textit{Memorisation}: $\tau_{\mathrm{train}}$ detected, $\taug$ not,
and $\tauf$ not detected.
\textit{Confusion}: $\tau_{\mathrm{train}}$ not detected (learning
failure; neither training nor test accuracy reaches the threshold).
\textit{Comprehension}: both $\tau_{\mathrm{train}}$ and $\taug$
detected with $\taug - \tau_{\mathrm{train}} < 2{,}000$ steps
(immediate generalisation; expected to be rare under this configuration).
Runs where $\tau_{\mathrm{train}}$ is detected, $\taug$ is not, but
$\tauf$ \emph{is} detected are labelled \textbf{F\_only} (Fourier
geometry without generalisation).
Each Grokking run is then labelled:
\textbf{G$<$F} ($\taug < \tauf$);
\textbf{F$<$G} ($\taug > \tauf$);
or \textbf{coincident} ($|\Dtau| \leq 500$ steps).
When a hyperparameter cell yields mixed phases across seeds (e.g.,
one seed Grokking and two seeds F\_only), the cell is reported as
\textbf{Mixed} with per-seed breakdowns shown in the table.

\subsection{Experimental Design}

The project will proceed in four stages.

\textbf{Stage~0 (Metric Validation).}
Six representative cells (two each from expected Grokking,
Memorisation, and Confusion regions) will be run with 3 seeds each
to validate that $\taug$ and $\tauf$ behave as expected and to
confirm that the chosen metric thresholds are appropriate before
committing to full-scale sweeps.

\textbf{Stage~1 (Coarse Grid).}
A 5$\times$5 grid spanning the ($\mathrm{lr}$, $\mathrm{wd}$)
hyperparameter space will be run with 1 seed per cell to map phase
boundaries and identify G$<$F signal.

\textbf{Stage~2 (Weight Decay Sweep).}
Fixed $\mathrm{lr} = 1.6\times10^{-3}$.
Weight decay: 10 log-spaced values in $[1.2, 3.5]$:
$\{1.20, 1.35, 1.52, 1.71, 1.93, 2.18, 2.45, 2.76, 3.11, 3.50\}$.
Seeds $\{42, 7, 2025\}$. Total: 30 runs.

\textbf{Stage~3 (Learning Rate Sweep).}
Fixed $\mathrm{wd} = 2.5$.
Learning rate: 10 log-spaced values in $[5\times10^{-4}, 8\times10^{-3}]$:
$\{5.0, 6.8, 9.3, 12.6, 17.1, 23.3, 31.7, 43.2, 58.8, 80.0\}
\times 10^{-4}$.
Seeds $\{42, 7, 2025\}$. Total: 30 runs.

\textbf{Sensitivity Analysis.}
Post-hoc: representative cells spanning strong, medium, and weak
G$<$F cases, any coincident cases, and any stable F$<$G cases
(if observed), will be re-analysed under all $3\times4\times3 = 36$
detector configurations (varying the null-update interval, EMA
smoothing coefficient $\alpha$, and BIC slope threshold).
Cells are selected by the following deterministic rule applied to
the Stage~2/3 outputs: rank all Grokking-phase cells by $|\Dtau|$
and pick the top, median, and bottom cells as the strong/medium/weak
G$<$F representatives; additionally include every cell with
$|\Dtau|\leq 500$ steps (coincident) and every cell with
$\Dtau<0$ (F$<$G), capped at three cells per category.
This yields a fixed, pre-registered cell set independent of which
configuration is used to compute $\Dtau$.
Figure~\ref{fig:sens_robustness} reports the headline robustness fraction
across configurations, Figure~\ref{fig:sens_percell} expands this into
per-cell $\Dtau$ breakdowns across the configuration grid, and
Figure~\ref{fig:sens_outlier} traces a single outlier cell to
illustrate when the changepoint is or is not accepted.

\begin{figure}[htbp]
  \centering
  \includegraphics[width=0.78\textwidth]{sensitivity_robustness_heatmap.pdf}
  \caption[Sensitivity robustness heatmap.]{Detector robustness across the $3\times4\times3$ configuration grid: each cell records the fraction of detector configurations under which the same G$<$F / F$<$G ordering is recovered.}
  \label{fig:sens_robustness}
\end{figure}

\begin{figure}[htbp]
  \centering
  \includegraphics[width=0.78\textwidth]{sensitivity_percell_heatmap.pdf}
  \caption[Per-cell sensitivity breakdown.]{Per-cell breakdown of $\Dtau$ across detector configurations, showing how the sign and magnitude of the lag respond to changes in the null-update interval, EMA $\alpha$, and BIC slope threshold.}
  \label{fig:sens_percell}
\end{figure}

\begin{figure}[htbp]
  \centering
  \includegraphics[width=0.78\textwidth]{sensitivity_outlier_trajectory.pdf}
  \caption[Outlier trajectory under varying detectors.]{Trajectory of $\fcorr(t)$ for a representative outlier cell under each detector configuration, illustrating the configurations under which the changepoint is or is not accepted.}
  \label{fig:sens_outlier}
\end{figure}

\section{Expected Outcomes}

\paragraph{Decision rule.}
We pre-register the following criterion for declaring the G$<$F
hypothesis \emph{supported} on a given sweep
(Stage~2 or Stage~3 considered separately):
(i)~at least 80\% of the Grokking-phase runs in that sweep have
$\Dtau > 0$ under the default detector configuration;
(ii)~the median of $\{\Dtau_i\}$ across those runs has a one-sided
$95\%$ bootstrap lower confidence bound \cite{efron1993bootstrap}
strictly greater than zero ($10{,}000$ resamples); and
(iii)~the sign of the median $\Dtau$ is preserved in
$\geq 30/36$ detector configurations from the sensitivity sweep.
F$<$G is declared supported by the symmetric criterion
($\Dtau<0$, upper CI bound below zero, sign preserved in
$\geq 30/36$ configurations).
Any other outcome---including a split between Stage~2 and Stage~3---is
reported as \emph{inconclusive}, and the per-axis breakdown is
presented as the primary contribution.

If the G$<$F hypothesis holds, we expect:
\begin{itemize}
  \item The majority of Grokking-phase runs will show positive $\Dtau$
    (generalisation precedes Fourier embedding alignment), with a
    consistent margin of several thousand training steps.
  \item The G$<$F pattern will persist across both the weight-decay axis
    (Stage~2) and the learning-rate axis (Stage~3), indicating it is a
    property of the Grokking phase rather than an artefact of a single
    hyperparameter setting.
  \item The ordering will be robust to detector hyperparameter choices:
    varying the null-update interval, EMA smoothing coefficient, and BIC
    slope threshold should not reverse the ordering for runs with large
    $|\Dtau|$.
\end{itemize}

The intended reporting format for the lag distribution and the per-run
$(\taug,\tauf)$ relationship is shown in Figures~\ref{fig:delta_hist}
and~\ref{fig:tau_scatter}.

\begin{figure}[htbp]
  \centering
  \includegraphics[width=0.78\textwidth]{delta_histogram.pdf}
  \caption[Distribution of $\Dtau$ across runs.]{Distribution of the signed lag $\Dtau = \tauf - \taug$ across Grokking-phase runs in the planned sweeps. Positive values correspond to G$<$F, negative to F$<$G; the bulk of the distribution is the primary test of the G$<$F hypothesis.}
  \label{fig:delta_hist}
\end{figure}

\begin{figure}[htbp]
  \centering
  \includegraphics[width=0.78\textwidth]{tau_scatter.pdf}
  \caption[Scatter of $\tauf$ versus $\taug$.]{Scatter of $\tauf$ against $\taug$ across runs. Points above the diagonal correspond to G$<$F orderings; points below to F$<$G.}
  \label{fig:tau_scatter}
\end{figure}

If F$<$G holds instead, or if the ordering is inconsistent across the
phase diagram, this would suggest that embedding-level Fourier structure
is a prerequisite for generalisation rather than a lagging consolidation
step---a qualitatively different picture of the grokking process.

A secondary expected outcome is a characterisation of \textbf{F\_only}
runs (Fourier geometry detected without generalisation), which would
illuminate whether embedding alignment is necessary, sufficient, or
neither for generalisation.

\section{Anticipated Limitations}

\begin{itemize}
  \item \textbf{Single task}: experiments are restricted to
    $(a+b)\bmod 53$; the generality of the G$<$F finding to other
    modular or non-modular tasks remains an open question.
  \item \textbf{1D confirmatory sweeps}: after the initial coarse 2D
    grid (Stage~1), each confirmatory stage fixes one hyperparameter
    axis; a dense 2D grid over both axes simultaneously would provide
    stronger coverage of the full Grokking region.
  \item \textbf{Resolution floor}: $\Dtau$ is measured at 500-step
    intervals; small lags ($|\Dtau|\lesssim500$ steps) cannot be
    reliably classified as one ordering or the other.
  \item \textbf{Circuit formation not measured}: this project measures
    the timing gap between generalisation onset and embedding-geometry
    onset; the earlier circuit-formation stage is not directly tracked
    in the same runs.
\end{itemize}

\section{Risk Assessment}

\subsection{OHS Risk Assessment}

This project is entirely computational: all work consists of writing and running Python code on a personal workstation, with no laboratory, field, or wet-lab work of any kind.
No hazardous materials, chemicals, biological agents, high-voltage
equipment, or physical prototypes are involved.
The project therefore falls under the \textbf{Low Risk} category
and is covered by the general OHS laboratory rules of the School
of EECS.
The main ergonomic consideration is prolonged computer use, which will be managed by taking regular breaks in line with standard university guidance.

\subsection{Ethics Assessment}

This project does not involve human participants, animal subjects, or the collection of personal data.
All experiments are conducted on synthetically generated arithmetic
datasets ($\{(a,b,c) \mid a+b \equiv c \pmod{53}\}$) with no
identifiable individuals.
No UQ ethical clearance is required.

\subsection{Project Schedule Risks}

\begin{itemize}
  \item \textbf{Computational runtime}: each full sweep (30 runs
    $\times$ 50{,}000--80{,}000 steps) may exceed initial time
    estimates if convergence is slow.
    \textit{Mitigation}: runs will be parallelised on a local GPU
    (RTX 4080).
    A pre-defined downgrade ladder will be applied if cumulative GPU
    time exceeds the budget by more than 20\%:
    (1)~drop Stage~3 seeds from $\{42, 7, 2025\}$ to $\{42, 2025\}$;
    (2)~halve the Stage~3 grid to 5 log-spaced points;
    (3)~cap the maximum step budget at 50{,}000 (no extension to
    80{,}000); (4)~as a last resort, restrict the sensitivity sweep
    to the 12-config subset corresponding to a single value of the
    BIC slope threshold.
    Each step preserves the primary G$<$F test on Stage~2 while
    progressively shedding lower-priority resolution.
  \item \textbf{Hardware failure}: GPU malfunction could interrupt
    experiments.
    \textit{Mitigation}: checkpoints will be saved every 500 steps; runs
    can be resumed or migrated to a cloud GPU instance.
  \item \textbf{Detector non-detection}: $\tauf$ may not be detected
    in a substantial fraction of runs, reducing sample size for the
    ordering analysis.
    \textit{Mitigation}: the F\_only label already accounts for this
    case; $\taug$ alone constitutes a valid secondary contribution,
    and the sensitivity sweep additionally reports the per-cell
    detection rate so non-detections are themselves an outcome.
\end{itemize}

\section{Project Milestones and Timeline}

The full milestone list, GPU time estimates, and current status are
summarised in Table~\ref{tab:milestones}.

\begin{table}[htbp]
  \centering
  \caption{Project milestones, planned timeline, and current status.}
  \label{tab:milestones}
  \renewcommand{\arraystretch}{1.2}
  \begin{tabular}{lllll}
    \toprule
    Milestone & Description & GPU est. & Planned Weeks & Status \\
    \midrule
    Stage~0 & Metric validation (18 runs) & $\sim$5 h & Weeks 1--2 & Complete \\
    Stage~1 & Coarse 5$\times$5 grid (25 runs) & $\sim$7 h & Weeks 3--4 & Complete \\
    Stage~2 & Weight-decay sweep (30 runs) & $\sim$9 h & Weeks 5--6 & Complete \\
    Stage~3 & Learning-rate sweep (30 runs) & $\sim$9 h & Weeks 7--8 & Complete \\
    Sensitivity & 36-config re-analysis (checkpoints) & $\sim$3 h & Weeks 9--10 & Complete \\
    Write-up  & Final thesis report & --- & Weeks 11--13 & In progress \\
    \bottomrule
  \end{tabular}
\end{table}

\subsection{Progress to date}

As of the time of submission, Stages~0 through~3 and the sensitivity
analysis have been completed on schedule.
Experiment data are archived under \texttt{runs/}.
The remaining deliverable is the final thesis write-up (Weeks~11--13),
which is currently in progress.
All experiments were run on a local GPU workstation
(NVIDIA RTX 4080 Laptop, 12~GB VRAM) using
Python~3.14, PyTorch~2.5, and CUDA~12.4.

{\renewcommand{\bibname}{References}%
\begin{thebibliography}{99}
\addcontentsline{toc}{chapter}{References}

\bibitem{nanda2023progress}
N.~Nanda, L.~Chan, T.~Lieberum, J.~Smith, and J.~Steinhardt,
``Progress measures for grokking via mechanistic interpretability,''
in \textit{Proc. Int. Conf. Learn. Representations (ICLR)}, 2023.

\bibitem{power2022grokking}
A.~Power, Y.~Burda, H.~Edwards, I.~Babuschkin, and V.~Misra,
``Grokking: Generalization beyond overfitting on small algorithmic
datasets,'' \textit{arXiv preprint arXiv:2201.02177}, 2022.

\bibitem{olah2020circuits}
C.~Olah, N.~Cammarata, L.~Schubert, G.~Goh, M.~Petrov, and
S.~Carter,
``Zoom in: An introduction to circuits,''
\textit{Distill}, vol.~5, no.~3, 2020, doi: 10.23915/distill.00024.001.

\bibitem{liu2022towards}
Z.~Liu, O.~Kitouni, N.~Nolte, E.~J.~Michaud, M.~Tegmark, and
M.~Williams, ``Towards understanding grokking: An effective theory
of representation learning,'' in \textit{Proc. Adv. Neural Inf.
Process. Syst. (NeurIPS)}, vol.~35, 2022, pp.~34651--34663.

\bibitem{park2024linear}
K.~Park, Y.~J.~Choe, and V.~Veitch,
``The linear representation hypothesis and the geometry of
large language models,''
in \textit{Proc. 41st Int. Conf. Mach. Learn. (ICML)},
\textit{Proc. Mach. Learn. Res.}, vol.~235, 2024,
pp.~39643--39666.

\bibitem{papyan2020prevalence}
V.~Papyan, X.~Y.~Han, and D.~L.~Donoho,
``Prevalence of neural collapse during the terminal phase of
deep learning training,''
\textit{Proc. Natl. Acad. Sci. U.S.A.}, vol.~117, no.~40,
pp.~24652--24663, 2020.

\bibitem{krogh1991weightdecay}
A.~Krogh and J.~A.~Hertz,
``A simple weight decay can improve generalization,''
in \textit{Proc. Adv. Neural Inf. Process. Syst. (NeurIPS)},
vol.~4, 1991, pp.~950--957.

\bibitem{thilak2022slingshot}
V.~Thilak, E.~Littwin, S.~Zhai, O.~Saremi, R.~Paiss, and
J.~Susskind,
``The slingshot mechanism: An empirical study of adaptive
optimizers and the grokking phenomenon,''
\textit{arXiv preprint arXiv:2206.04817}, 2022.

\bibitem{notsawo2023predicting}
P.~J.~T.~Notsawo, H.~Zhou, M.~Pezeshki, I.~Rish, and G.~Dumas,
``Predicting grokking long before it happens: A look into the
loss landscape of models which grok,''
\textit{arXiv preprint arXiv:2306.13253}, 2023.

\bibitem{davies2023unifying}
X.~Davies, L.~Langosco, and D.~Krueger,
``Unifying grokking and double descent,''
\textit{arXiv preprint arXiv:2303.06173}, 2023.

\bibitem{merrill2023tale}
W.~Merrill, N.~Tsilivis, and A.~Shukla,
``A tale of two circuits: Grokking as competition of sparse and
dense subnetworks,''
\textit{arXiv preprint arXiv:2303.11873}, 2023.

\bibitem{liu2023omnigrok}
Z.~Liu, E.~J.~Michaud, and M.~Tegmark,
``OmniGrok: Grokking beyond algorithmic data,''
in \textit{Proc. Int. Conf. Learn. Representations (ICLR)}, 2023.

\bibitem{lyu2024dichotomy}
K.~Lyu, J.~Jin, Z.~Li, S.~S.~Du, J.~D.~Lee, and W.~Hu,
``Dichotomy of early and late phase implicit biases can provably
induce grokking,''
in \textit{Proc. Int. Conf. Learn. Representations (ICLR)}, 2024.

\bibitem{varma2023explaining}
V.~Varma, R.~Shah, Z.~Kenton, J.~Kram\'{a}r, and R.~Kumar,
``Explaining grokking through circuit efficiency,''
\textit{arXiv preprint arXiv:2309.02390}, 2023.

\bibitem{zhong2023clock}
Z.~Zhong, Z.~Liu, M.~Tegmark, and J.~Andreas,
``The clock and the pizza: Two stories in mechanistic explanation
of neural networks,''
in \textit{Proc. Adv. Neural Inf. Process. Syst. (NeurIPS)},
vol.~36, 2023.

\bibitem{gromov2023grokking}
A.~Gromov,
``Grokking modular arithmetic,''
\textit{arXiv preprint arXiv:2301.02679}, 2023.

\bibitem{saxe2014exact}
A.~M.~Saxe, J.~L.~McClelland, and S.~Ganguli,
``Exact solutions to the nonlinear dynamics of learning in deep
linear neural networks,''
in \textit{Proc. Int. Conf. Learn. Representations (ICLR)}, 2014.

\bibitem{barak2022hidden}
B.~Barak, B.~L.~Edelman, S.~Goel, S.~Kakade, E.~Malach, and
C.~Zhang, ``Hidden progress in deep learning: SGD learns parities
near the computational limit,'' in \textit{Proc. Adv. Neural Inf.
Process. Syst. (NeurIPS)}, vol.~35, 2022, pp.~21750--21764.

\bibitem{elhage2021mathematical}
N.~Elhage \textit{et al.},
``A mathematical framework for transformer circuits,''
\textit{Transformer Circuits Thread}, 2021.
[Online]. Available:
\url{https://transformer-circuits.pub/2021/framework/index.html}.
Accessed: Apr.\ 2026.

\bibitem{rubin2024grokking}
N.~Rubin, I.~Seroussi, and Z.~Ringel,
``Grokking as a first-order phase transition in two-layer networks,''
in \textit{Proc. Int. Conf. Learn. Representations (ICLR)}, 2024.

\bibitem{vaswani2017attention}
A.~Vaswani, N.~Shazeer, N.~Parmar, J.~Uszkoreit, L.~Jones,
A.~N.~Gomez, {\L}.~Kaiser, and I.~Polosukhin,
``Attention is all you need,''
in \textit{Proc. Adv. Neural Inf. Process. Syst. (NeurIPS)},
vol.~30, 2017, pp.~5998--6008.

\bibitem{radford2018gpt}
A.~Radford, K.~Narasimhan, T.~Salimans, and I.~Sutskever,
``Improving language understanding by generative pre-training,''
OpenAI, Tech. Rep., 2018.
[Online]. Available:
\url{https://cdn.openai.com/research-covers/language-unsupervised/language_understanding_paper.pdf}.
Accessed: Apr.\ 2026.

\bibitem{he2016resnet}
K.~He, X.~Zhang, S.~Ren, and J.~Sun,
``Deep residual learning for image recognition,''
in \textit{Proc. IEEE Conf. Comput. Vis. Pattern Recognit. (CVPR)},
2016, pp.~770--778.

\bibitem{ba2016layernorm}
J.~L.~Ba, J.~R.~Kiros, and G.~E.~Hinton,
``Layer normalization,''
\textit{arXiv preprint arXiv:1607.06450}, 2016.

\bibitem{loshchilov2019adamw}
I.~Loshchilov and F.~Hutter,
``Decoupled weight decay regularization,''
in \textit{Proc. Int. Conf. Learn. Representations (ICLR)}, 2019.

\bibitem{kingma2015adam}
D.~P.~Kingma and J.~Ba,
``Adam: A method for stochastic optimization,''
in \textit{Proc. Int. Conf. Learn. Representations (ICLR)}, 2015.

\bibitem{schwarz1978bic}
G.~Schwarz,
``Estimating the dimension of a model,''
\textit{Ann. Statist.}, vol.~6, no.~2, pp.~461--464, 1978.

\bibitem{muggeo2003segmented}
V.~M.~R.~Muggeo,
``Estimating regression models with unknown break-points,''
\textit{Statist. Med.}, vol.~22, no.~19, pp.~3055--3071, 2003.

\bibitem{efron1993bootstrap}
B.~Efron and R.~J.~Tibshirani,
\textit{An Introduction to the Bootstrap}.
New York, NY, USA: Chapman and Hall, 1993.

\end{thebibliography}

\end{document}
